\section{Traitement des erreurs}\label{sec:traitement-des-erreurs}

\subsection{Introduction}\label{subsec:introduction}
Cette section présente les principaux choix de réalisation concernant le traitement des erreurs dans le projet MiniJAJA. Le traitement des erreurs est un aspect crucial du compilateur et de l'interpréteur, permettant de fournir des retours clairs et précis aux utilisateurs lors de la détection d'anomalies.

\subsection{Stratégie générale de gestion des erreurs}\label{subsec:strategie-generale-de-gestion-des-erreurs}

\subsubsection{Classification des erreurs}
Le projet MiniJAJA gère plusieurs types d'erreurs à travers un système de diagnostic unifié.
Le système distingue quatre phases d'erreurs, définies dans l'énumération \texttt{Phase} :
\begin{itemize}
    \item \textbf{Erreurs lexicales} (\texttt{LEXICAL}) : détectées lors de l'analyse lexicale (tokens invalides, caractères non reconnus)
    \item \textbf{Erreurs syntaxiques} (\texttt{SYNTAX}) : détectées lors de l'analyse syntaxique (règles de grammaire violées)
    \item \textbf{Erreurs sémantiques} (\texttt{SEMANTIC}) : détectées lors de l'analyse sémantique (typage, portée, déclarations)
    \item \textbf{Erreurs d'exécution} (\texttt{RUNTIME}) : détectées lors de l'interprétation du code JajaCode
\end{itemize}

Chaque erreur possède également une sévérité définie par l'énumération \texttt{Severity} :
\begin{itemize}
    \item \textbf{ERROR} : erreur bloquante empêchant la compilation ou l'exécution
    \item \textbf{WARNING} : avertissement n'empêchant pas la compilation, mais signalant un problème potentiel.
    \textit{Note : cette fonctionnalité est définie dans l'architecture, mais n'est pas encore exploitée dans l'implémentation actuelle. L'infrastructure est en place pour permettre une utilisation future.}
\end{itemize}

\subsubsection{Architecture du système de diagnostic}

Le système de gestion des erreurs repose sur une architecture modulaire composée de plusieurs classes clés :

\paragraph{Classe Diagnostic}
La classe \texttt{Diagnostic} est un record Java qui encapsule toutes les informations d'une erreur :

\vspace{0.3cm}
\begin{lstlisting}[language=java, caption=Structure du Diagnostic,label={lst:diagnostic-record}]
public record Diagnostic(
    Severity severity,      // ERROR ou WARNING
    Phase phase,            // LEXICAL, SYNTAX, SEMANTIC ou RUNTIME
    SourcePosition position, // Fichier, ligne, colonne
    String message          // Message descriptif
) {}
\end{lstlisting}

\paragraph{Classe SourcePosition}
La classe \texttt{SourcePosition} capture la position exacte de l'erreur dans le code source, incluant le nom du fichier, le numéro de ligne et la position dans la ligne.
Cette information permet aux développeurs de localiser rapidement l'origine du problème.

\paragraph{Classe DiagnosticCollector}
Le \texttt{DiagnosticCollector} est responsable de collecter tous les diagnostics durant les phases d'analyse.
Il permet de :
\begin{itemize}
    \item Accumuler plusieurs erreurs sans interrompre l'analyse
    \item Formater tous les diagnostics en une seule chaîne de caractères via \texttt{formatDiagnostics()}
    \item Vérifier la présence d'erreurs spécifiques par phase :
    \begin{itemize}
        \item \texttt{hasErrors()} : présence d'erreurs de type ERROR
        \item \texttt{hasSyntaxErrors()} : erreurs syntaxiques ou lexicales (phases SYNTAX ou LEXICAL)
        \item \texttt{hasSemanticErrors()} : erreurs sémantiques (phase SEMANTIC)
    \end{itemize}
\end{itemize}

\subsubsection{Principes de conception}
Les choix de conception pour le traitement des erreurs suivent plusieurs principes :
\begin{itemize}
    \item \textbf{Clarté} : les messages d'erreur sont explicites et compréhensibles
    \item \textbf{Précision} : indication exacte de la position (fichier, ligne, colonne)
    \item \textbf{Récupération} : poursuite de l'analyse pour détecter plusieurs erreurs simultanément
    \item \textbf{Cohérence} : uniformité du format des messages via le système \texttt{Diagnostic}
    \item \textbf{Traçabilité} : chaque erreur est associée à sa phase de détection
\end{itemize}

\subsection{Traitement des erreurs par module}\label{subsec:traitement-des-erreurs-par-module}

\subsubsection{Module LexerParser}

\paragraph{Erreurs lexicales et syntaxiques}

Le module LexerParser utilise ANTLR4 pour l'analyse lexicale et syntaxique.
Les erreurs sont capturées via un listener personnalisé :

\begin{lstlisting}[language=java, caption=SyntaxErrorListener,label={lst:syntaxerrorlistener}]
public class SyntaxErrorListener extends BaseErrorListener {
    private final DiagnosticCollector diagnosticCollector;
    private final String fileName;

    @Override
    public void syntaxError(Recognizer<?, ?> recognizer,
                           Object offendingSymbol,
                           int line, int charPositionInLine,
                           String msg, RecognitionException e) {
        SourcePosition pos = new SourcePosition(fileName,
                                                line,
                                                charPositionInLine);
        diagnosticCollector.report(Severity.ERROR,
                                  Phase.SYNTAX,
                                  pos, msg);
    }
}
\end{lstlisting}

Le \texttt{SyntaxErrorListener} :
\begin{itemize}
    \item Remplace les listeners d'erreur par défaut d'ANTLR
    \item Capture les erreurs lexicales et syntaxiques
    \item Les transforme en objets \texttt{Diagnostic} structurés
    \item Permet la collecte de multiples erreurs sans interruption
\end{itemize}

\paragraph{Erreurs sémantiques}

L'analyse sémantique est réalisée par le \texttt{MiniJajaSemanticAnalyzer}, qui délègue les vérifications à des composants spécialisés :

\begin{itemize}
    \item \textbf{TypeChecker} : vérifie la compatibilité des types dans les expressions et affectations
    \item \textbf{ScopeResolver} : résout les références aux variables et vérifie leur portée
    \item \textbf{DeclarationCollector} : collecte les déclarations et détecte les redéclarations
    \item \textbf{MethodCallValidator} : valide les appels de méthodes (existence, arité, types des paramètres)
    \item \textbf{TypeInferenceEngine} : infère les types des expressions complexes
\end{itemize}

Les erreurs sémantiques typiques incluent :
\begin{itemize}
    \item Variables non déclarées
    \item Redéclaration de variables
    \item Incompatibilités de types (ex: affectation d'un entier à un booléen)
    \item Appels de méthodes inexistantes
    \item Mauvais nombre d'arguments dans les appels de méthode
    \item Accès à des tableaux avec un indice non entier
    \item Utilisation de variables hors de leur portée
\end{itemize}

Chaque validateur utilise un \texttt{SemanticContext} partagé qui maintient :
\begin{itemize}
    \item La table des symboles (variables, méthodes, constantes)
    \item Les informations de typage
    \item La portée courante
\end{itemize}

\subsubsection{Interpréteur JajaCode}

L'interpréteur JajaCode utilise une hiérarchie d'exceptions spécifiques pour les erreurs d'exécution, toutes héritant de \texttt{JajaCodeRuntimeException}.

\paragraph{Exception de base : JajaCodeRuntimeException}

\vspace{0.3cm}
\begin{lstlisting}[language=java, caption=JajaCodeRuntimeException,label={lst:jjcruntimeexception}]
public class JajaCodeRuntimeException extends RuntimeException {
    private final int programCounter;
    private final String axiomeName;

    public JajaCodeRuntimeException(String message,
                                    String axiomeName,
                                    int programCounter) {
        super(String.format("[PC=%d, Axiome=%s] %s",
                           programCounter, axiomeName, message));
        this.axiomeName = axiomeName;
        this.programCounter = programCounter;
    }
}
\end{lstlisting}

Cette exception de base fournit :
\begin{itemize}
    \item Le \textbf{Program Counter (PC)} : position dans le code JajaCode où l'erreur s'est produite.
    \item Le \textbf{nom de l'axiome} : l'instruction JajaCode en cours d'exécution.
    \item Un \textbf{message formaté} incluant ces informations contextuelles.
\end{itemize}

\paragraph{Exceptions spécifiques}

Le projet définit plusieurs exceptions spécialisées :

\begin{itemize}
    \item \textbf{DivisionByZeroException} : division ou modulo par zéro lors d'opérations arithmétiques,
    \item \textbf{StackUnderflowException} : tentative de pop sur une pile vide,
    \item \textbf{TypeMismatchException} : incompatibilité de types à l'exécution,
    \item \textbf{InvalidAddressException} : accès mémoire à une adresse invalide,
    \item \textbf{AssignmentException} : erreur lors d'une affectation,
    \item \textbf{UndefinedSymbolException} : référence à un symbole non défini.
\end{itemize}

Exemple d'exception spécialisée :
\begin{lstlisting}[language=java, caption=DivisionByZeroException,label={lst:divisionbyzero-exception}]
public class DivisionByZeroException
    extends JajaCodeRuntimeException {

    public DivisionByZeroException(String axiomeName,
                                   int programCounter) {
        super("Division par zero", axiomeName, programCounter);
    }
}
\end{lstlisting}

\subsection{Format des messages d'erreur}\label{subsec:format-des-messages-d'erreur}

\subsubsection{Structure des messages}

Les messages d'erreur suivent un format uniforme qui dépend de la phase :

\paragraph{Erreurs de compilation (Lexicales, Syntaxiques, Sémantiques)}

\vspace{0.3cm}
\begin{lstlisting}[caption=Format des erreurs de compilation,label={lst:ex-err-compil}]
[PHASE] fichier.mjj:ligne:colonne - Message descriptif
\end{lstlisting}

Exemple :
\begin{lstlisting}[caption=Exemple d'erreur sémantique,label={lst:ex-err-semantique}]
[SEMANTIC] test.mjj:15:8 - Variable 'x' non declaree
\end{lstlisting}

\paragraph{Erreurs d'exécution (Runtime JajaCode)}

\vspace{0.3cm}
\begin{lstlisting}[caption=Format des erreurs d'execution,label={lst:ex-err-runtime}]
[PC=<compteur>, Axiome=<nom>] Message descriptif
\end{lstlisting}

Exemple :
\begin{lstlisting}[caption=Exemple d'erreur d'execution,label={lst:ex-err-execution}]
[PC=42, Axiome=div] Division par zero
\end{lstlisting}

\subsubsection{Informations contextuelles}

Chaque message d'erreur inclut des informations contextuelles permettant au développeur de localiser et comprendre le problème :
\begin{itemize}
    \item \textbf{Nom du fichier} : pour les projets multi-fichiers
    \item \textbf{Ligne et colonne} : position exacte dans le source
    \item \textbf{Program Counter} : pour les erreurs JajaCode, position dans le bytecode
    \item \textbf{Nom de l'axiome} : instruction JajaCode en cours d'exécution
    \item \textbf{Message descriptif} : explication claire du problème
\end{itemize}

\subsection{Gestion des exceptions}\label{subsec:gestion-des-exceptions}

\subsubsection{Hiérarchie des exceptions}

Le projet utilise une hiérarchie d'exceptions bien structurée organisée en trois catégories principales :

\begin{lstlisting}[caption=Hierarchie des exceptions,label={lst:hierarchie-des-exceptions}]
java.lang.RuntimeException
  |
  +-- SyntaxException (erreurs syntaxiques et lexicales)
  |
  +-- SemanticException (erreurs semantiques)
  |
  +-- JajaCodeRuntimeException (erreurs d'execution)
        |
        +-- DivisionByZeroException
        +-- StackUnderflowException
        +-- TypeMismatchException
        +-- InvalidAddressException
        +-- AssignmentException
        +-- UndefinedSymbolException
\end{lstlisting}

Cette hiérarchie permet :
\begin{itemize}
    \item Une capture sélective des exceptions par phase de compilation
    \item Un traitement différencié selon le type d'erreur
    \item Une extension facile avec de nouvelles exceptions
    \item Une séparation claire entre erreurs de compilation et d'exécution
\end{itemize}

\paragraph{Exceptions de compilation}

Les exceptions \texttt{SyntaxException} et \texttt{SemanticException} sont définies dans le package \texttt{fr.ufrst.m1info.gl.groupe7.lexerparser.errors.exceptions} et suivent un pattern unifié :

\begin{lstlisting}[language=java, caption=Pattern des exceptions de compilation,label={lst:compilation-exception-pattern}]
public class SemanticException extends RuntimeException {
    public SemanticException(DiagnosticCollector collector) {
        super(collector != null && collector.hasErrors()
                ? collector.formatDiagnostics()
                : "Unknown semantic error.");
    }
}
\end{lstlisting}

Ce pattern présente plusieurs avantages :
\begin{itemize}
    \item \textbf{Intégration avec DiagnosticCollector} : le message d'exception contient automatiquement tous les diagnostics formatés
    \item \textbf{Messages riches} : inclusion de toutes les erreurs détectées, pas seulement la première
    \item \textbf{Cohérence} : format uniforme pour toutes les exceptions de compilation
    \item \textbf{Robustesse} : message par défaut si le collector est null ou vide
\end{itemize}

\texttt{SyntaxException} est levée lorsque des erreurs lexicales ou syntaxiques sont détectées (phase LEXICAL ou SYNTAX), tandis que \texttt{SemanticException} est levée pour les erreurs d'analyse sémantique (phase SEMANTIC).

\subsubsection{Propagation et capture}

Les exceptions sont propagées jusqu'au point d'entrée principal où elles sont capturées et transformées en messages utilisateur. Le système de \texttt{DiagnosticCollector} permet également de collecter plusieurs erreurs avant de les rapporter toutes ensemble.

\subsection{Tests et validation}\label{subsec:tests-et-validation}

\subsubsection{Couverture des tests}

Le module LexerParser contient \textbf{38 classes de tests}, dont plusieurs dédiées spécifiquement au traitement des erreurs :

\begin{itemize}
    \item \texttt{ExceptionTypeTest} : vérifie que les bonnes exceptions sont levées selon le type d'erreur (syntaxique vs sémantique)
    \item \texttt{SyntaxErrorListenerTest} : teste la capture des erreurs syntaxiques via le listener ANTLR
    \item \texttt{MiniJajaSemanticAnalyzerTest} : teste l'analyse sémantique générale
    \item \texttt{MiniJajaSemanticAnalyzerScopeTest} : teste la vérification de portée des variables
    \item \texttt{MiniJajaSemanticAnalyzerReturnTest} : teste la validation des instructions de retour
    \item \texttt{DiagnosticCollectorTest} : teste le collecteur de diagnostics et ses méthodes de filtrage
    \item \texttt{DiagnosticTest} : teste la structure des diagnostics (record)
    \item \texttt{JajaCodeExceptionsTest} : teste les exceptions d'exécution JajaCode
    \item \texttt{ArrayAssignmentTest} : teste les erreurs d'affectation de tableaux
    \item \texttt{TypeIncompatibilityTest} : teste les incompatibilités de types
\end{itemize}

\subsubsection{Stratégie de test}

Les tests vérifient :
\begin{itemize}
    \item La \textbf{détection correcte} des erreurs (pas de faux négatifs)
    \item Le \textbf{type d'exception} levé selon la phase d'erreur
    \item Le \textbf{format des messages} d'erreur intégrés aux exceptions
    \item La \textbf{position exacte} (ligne/colonne) des erreurs
    \item La \textbf{phase correcte} de détection
    \item La \textbf{récupération} après erreur (continuation de l'analyse)
\end{itemize}

Exemple de test vérifiant le type d'exception :
\begin{lstlisting}[language=java, caption=Test de type d'exception semantique,label={lst:test-semantic-exception}]
@Test
@DisplayName("Semantic error should throw SemanticException")
void semanticError_shouldThrowSemanticException() {
    String codeWithSemanticError = """
        class C {
            main {
                int x = 5;
                int x = 10;  // Redeclaration
            }
        }
    """;

    assertThrows(SemanticException.class,
        () -> MiniJajaCompiler
            .getMiniJajaCompilerVisitorFromString(codeWithSemanticError),
        "Semantic errors should throw SemanticException");
}
\end{lstlisting}

Ce test vérifie que la redéclaration d'une variable lève bien une \texttt{SemanticException} et non une \texttt{SyntaxException}. D'autres tests similaires vérifient les incompatibilités de types, les variables non déclarées, et les erreurs syntaxiques.

\subsection{Améliorations futures}\label{subsec:ameliorations-futures}

Plusieurs améliorations sont envisageables pour le système de gestion des erreurs :

\begin{itemize}
    \item \textbf{Activation des warnings} : exploiter le niveau de sévérité \texttt{WARNING} déjà présent dans l'architecture pour signaler des problèmes potentiels sans bloquer la compilation. Exemples d'usage :
    \begin{itemize}
        \item Variables déclarées mais jamais utilisées
        \item Code mort (unreachable code)
        \item Comparaisons redondantes
        \item Conventions de nommage non respectées
        \item Déprécations
    \end{itemize}
    \item \textbf{Suggestions de correction} : proposer des corrections automatiques (ex: ``Vouliez-vous dire 'integer' ?'')
    \item \textbf{Messages multilingues} : support de l'internationalisation des messages
    \item \textbf{Colorisation} : messages avec codes couleur dans la CLI pour une meilleure lisibilité
    \item \textbf{Extraits de code} : afficher l'extrait de code source avec l'erreur soulignée
    \item \textbf{Erreurs catégorisées} : codes d'erreur numérotés pour une documentation détaillée
    \item \textbf{Mode strict/permissif} : niveau de sévérité configurable (warnings as errors, ignorer certains warnings)
\end{itemize}

\subsection{Conclusion}\label{subsec:conclusion}

Le système de gestion des erreurs du projet MiniJAJA repose sur une architecture robuste et extensible.
L'utilisation d'un système de diagnostic unifié à travers toutes les phases (lexicale, syntaxique, sémantique, exécution) garantit la cohérence des messages d'erreur.

La hiérarchie d'exceptions bien structurée distingue clairement trois catégories d'erreurs :
\begin{itemize}
    \item \texttt{SyntaxException} pour les erreurs de parsing (lexicales et syntaxiques)
    \item \texttt{SemanticException} pour les erreurs d'analyse sémantique
    \item \texttt{JajaCodeRuntimeException} et ses sous-classes pour les erreurs d'exécution
\end{itemize}

L'intégration étroite entre \texttt{DiagnosticCollector} et les exceptions de compilation permet de fournir des messages d'erreur riches et détaillés, incluant tous les problèmes détectés plutôt que seulement le premier.
Les tests exhaustifs, incluant \texttt{ExceptionTypeTest} pour vérifier le type d'exception levé, assurent la fiabilité du système.

Cette approche permet aux utilisateurs de comprendre rapidement les problèmes dans leur code et de les corriger efficacement, tout en facilitant la maintenance et l'évolution future du système de gestion des erreurs.
